\documentclass[12pt]{article}
\usepackage[utf8]{inputenc}
\usepackage[T1]{fontenc}
\usepackage{hyperref}
\title{Balangingih Sama Language Resources for K-12 Curriculum Development}
\author{Compiled by Claude Sonnet 5.5 (claude-sonnet-5-5) for Andr\'as Kornai}
\date{}
\begin{document}
\maketitle

\section{Introduction}
Balangingih Sama (ISO 639-3 sse; Glottocode bala1311) is an Austronesian language of the Sama--Bajaw group. It is spoken mainly in the Philippines, with a smaller community in Malaysia. Alternate names include Balangingi, Bangingi, Northern Sama and Samal. Joshua Project lists dialects Balangingi, Daongdung, Kabinga'an, Lutangan and Sibuguey. Speaker estimates there are about 125,000--126,000 in the Philippines and about 6,000 in Malaysia. The task description gives about 100,000. Wikitongues reports both Arabic and Latin scripts in use.

\section{Primary Educational Resources}
No primers, textbooks, or school materials were found in search. No evidence of school use was found.

\section{Linguistic and Reference Materials}
\begin{itemize}
\item WALS entry; its reference source is Gault (1979) on the phonemics and morphophonemics of Sama Balangingi: \url{https://wals.info/languoid/lect/wals_code_sbg}
\item Wikitongues language page: \url{https://wikitongues.org/languages/sse/}
\end{itemize}

\section{Religious and Community Resources}
Joshua Project's language page reports a New Testament translation and describes the community as Muslim: \url{https://joshuaproject.net/languages/sse}. Its prayer card for the people group is at \url{https://joshuaproject.net/pray/cards/single/14031/RP}.

\section{Digital and Audio Resources}
No audio, video, or corpus resources were identified in search.

\section{Curriculum Development Context}
Very little is documented. Resources for neighbouring Sama varieties could be consulted, but should be kept separate from Balangingih materials. Local orthography decisions (Latin versus Arabic) would need community input.

\end{document}